\documentclass[11pt,a4paper]{article}

\usepackage[margin=1in]{geometry}
\usepackage[T1]{fontenc}
\usepackage{graphicx}
\usepackage{xcolor}
\usepackage{listings}
\usepackage{hyperref}
\usepackage{enumitem}
\usepackage{fancyvrb}

\hypersetup{
    colorlinks=true,
    linkcolor=black,
    urlcolor=blue
}

\definecolor{codebg}{RGB}{245,245,245}
\definecolor{codekw}{RGB}{0,90,160}
\definecolor{codecm}{RGB}{0,128,0}
\definecolor{codest}{RGB}{170,55,30}

\lstdefinestyle{reportcode}{
    backgroundcolor=\color{codebg},
    basicstyle=\ttfamily\small,
    keywordstyle=\color{codekw}\bfseries,
    commentstyle=\color{codecm}\itshape,
    stringstyle=\color{codest},
    breaklines=true,
    showstringspaces=false,
    frame=single,
    framerule=0.4pt,
    rulecolor=\color{black!30},
    xleftmargin=10pt,
    xrightmargin=5pt,
    aboveskip=10pt,
    belowskip=10pt
}

\lstdefinelanguage{ebnf}{
    morecomment=[l]{//},
    morestring=[b]",
    sensitive=true
}

\lstset{style=reportcode}

\newcommand{\imgplaceholder}[2][0.88\linewidth]{%
    \begin{center}
        \fbox{\parbox[c][0.26\textheight][c]{#1}{\centering\itshape #2}}
    \end{center}%
}

\newcommand{\screenshotplaceholder}[2]{%
    \begin{figure}[h!]
        \centering
        \imgplaceholder{Screenshot placeholder\\[4pt]Replace with image from:\\ #2}
        \caption{#1}
    \end{figure}
}

\begin{document}

\begin{titlepage}
    \centering
    \vspace*{2cm}

    {\scshape\LARGE Compiler Construction Project \par}
    \vspace{0.5cm}
    {\large Spring 2026\par}

    \vspace{2.2cm}
    \rule{\linewidth}{0.5pt}\\[0.6cm]
    {\huge\bfseries Compiler for a C-Like Language Subset\par}
    \vspace{0.3cm}
    {\Large\itshape VM-Based Code Generation and Runtime Execution\par}
    \rule{\linewidth}{0.5pt}\\[1.5cm]

    \vspace{1.2cm}
    {\large\bfseries Submitted by\par}
    \vspace{0.6cm}

    \begin{tabular}{ll}
        \textbf{Name} & \textbf{CMS ID} \\
        \rule{0pt}{14pt}\textit{Aqsa Batool} & \textit{413777} \\
        \rule{0pt}{14pt}\textit{Muhammad Hamza} & \textit{407251} \\
        \rule{0pt}{14pt}\textit{Ahmed Mohiuddin Shah} & \textit{415216}
    \end{tabular}

    \vfill

    {\large Project Report\par}
    {\large \today\par}
\end{titlepage}

\section{Introduction}

This project implements a compiler for a well-defined subset of the C language.
The implementation covers the complete frontend and execution pipeline: lexical
analysis, syntax analysis, semantic analysis, intermediate representation,
object-code generation, and runtime execution. Instead of generating native
machine code, the compiler emits a custom human-readable VM object format
(``.vmo''), and a separate runtime executes that object code. This design keeps
the compilation phases explicit and easy to demonstrate while still satisfying
the requirement that user programs can be compiled and executed.

The project is organized around two tools:

\begin{itemize}[leftmargin=2em]
    \item \textbf{cforge}: compiler executable that tokenizes, parses, checks,
    lowers, and emits VM object code.
    \item \textbf{crun}: runtime executable that loads the generated object file
    and executes it on a custom virtual machine.
\end{itemize}

The project also includes phase-wise inspection flags so that every important
stage can be shown independently during presentation or debugging.

\section{Objectives}

The main objectives achieved in this project are:

\begin{enumerate}[leftmargin=2em]
    \item Design a handwritten compiler for a meaningful subset of C.
    \item Implement lexical analysis with position-aware tokens.
    \item Implement recursive-descent parsing and explicit abstract syntax tree
    construction.
    \item Perform semantic analysis including scope handling, type checking,
    function checking, and array checking.
    \item Generate an intermediate representation and lower it to VM object
    code.
    \item Execute compiled programs through a separate runtime.
    \item Provide clean build, test, compile, and run workflows.
\end{enumerate}

\section{Implemented Language Subset}

The implemented subset is intentionally smaller than full ISO C, but large
enough to demonstrate all core compiler phases.

\subsection*{Types}

\begin{itemize}[leftmargin=2em]
    \item \texttt{int}
    \item \texttt{char}
    \item \texttt{float}
    \item \texttt{void}
\end{itemize}

\subsection*{Declarations}

\begin{itemize}[leftmargin=2em]
    \item global scalar declarations
    \item global 1D fixed-size array declarations
    \item local scalar declarations
    \item local 1D fixed-size array declarations
    \item scalar initializers
    \item \texttt{char s[N] = "literal"}
\end{itemize}

\subsection*{Statements}

\begin{itemize}[leftmargin=2em]
    \item block statements
    \item assignments
    \item expression statements
    \item \texttt{if} / \texttt{else}
    \item \texttt{while}
    \item \texttt{return}
\end{itemize}

\subsection*{Expressions}

\begin{itemize}[leftmargin=2em]
    \item identifiers
    \item integer, float, char, and string literals
    \item function calls
    \item array indexing
    \item unary \texttt{-} and \texttt{!}
    \item binary \texttt{+ - * / \% < <= > >= == != \&\& ||}
\end{itemize}

\subsection*{Functions}

\begin{itemize}[leftmargin=2em]
    \item user-defined functions
    \item typed parameters
    \item return values
    \item \texttt{main}
    \item function must be defined before call
\end{itemize}

\subsection*{Builtins}

\begin{itemize}[leftmargin=2em]
    \item \texttt{print\_int(int)}
    \item \texttt{print\_float(float)}
    \item \texttt{print\_char(char)}
    \item \texttt{print\_str(char[])}
    \item \texttt{print\_newline()}
\end{itemize}

\subsection*{Current Deliberate Limits}

\begin{itemize}[leftmargin=2em]
    \item no \texttt{for} statement
    \item no 2D arrays
    \item no array parameters
    \item no optimizer pass yet
    \item no escape sequences in char and string literals yet
    \item \texttt{\&\&} and \texttt{||} evaluate both operands
\end{itemize}

\section{Project Structure}

\begin{lstlisting}
compiler/
|-- cmd/
|   |-- cforge/          # compiler entry point
|   `-- crun/            # runtime entry point
|-- internal/
|   |-- token/           # token kinds and positions
|   |-- lexer/           # handwritten scanner
|   |-- ast/             # AST node definitions and pretty printer
|   |-- parser/          # recursive-descent parser
|   |-- sema/            # semantic analysis and symbol checking
|   |-- ir/              # intermediate representation and lowering
|   |-- vm/              # VM object model, loader, and executor
|   |-- driver/          # compiler CLI logic
|   `-- runtimecli/      # runtime CLI logic
|-- samples/
|   |-- pass/            # valid sample programs
|   `-- fail/            # invalid sample programs
|-- Makefile
|-- README.md
`-- report.tex
\end{lstlisting}

\section{Compiler Pipeline Overview}

The implemented compilation and execution flow is:

\begin{lstlisting}
source.c
  -> lexer
  -> tokens
  -> parser
  -> AST
  -> semantic analysis
  -> IR
  -> VM object code (.vmo)
  -> crun runtime
  -> observed program output
\end{lstlisting}

This separation is useful for both grading and debugging because each stage can
be displayed independently using compiler flags.

\section{Lexical Analysis}

The lexer is handwritten in Go. It reads the source character stream and emits a
sequence of tokens with line and column positions. These positions are reused by
the parser and semantic analyzer to report precise diagnostics.

The scanner recognizes:

\begin{itemize}[leftmargin=2em]
    \item keywords: \texttt{int}, \texttt{char}, \texttt{float}, \texttt{void},
    \texttt{if}, \texttt{else}, \texttt{while}, \texttt{return}
    \item identifiers
    \item integer literals
    \item float literals
    \item character literals
    \item string literals
    \item operators and punctuation
    \item single-line comments beginning with \texttt{//}
\end{itemize}

The compiler intentionally supports only plain character and string literals in
this version, with no escape-sequence processing.

Example token dump produced by:

\begin{lstlisting}
./bin/cforge --tokens samples/slice7_run.c
\end{lstlisting}

\begin{lstlisting}
1:1 KW_INT       "int"
2:5 KW_CHAR      "char"
2:14 CHAR_LIT     "'A'"
5:15 STRING_LIT   "\"hello\""
9:1 EOF          ""
\end{lstlisting}

\section{Grammar Specification}

The grammar is implemented with handwritten recursive-descent parsing. The
following EBNF describes the accepted subset.

\subsection*{Top-Level Declarations}

\begin{lstlisting}[language=ebnf]
program        = { top_decl } EOF ;

top_decl       = var_decl | function_def ;

function_def   = type_spec IDENTIFIER "(" params ")" block ;

params         = "void"
               | param { "," param }
               | empty ;

param          = type_spec IDENTIFIER ;

var_decl       = type_spec IDENTIFIER [ array_spec ] [ "=" expr ] ";" ;

array_spec     = "[" INT_LITERAL "]" ;

type_spec      = "int" | "char" | "float" | "void" ;
\end{lstlisting}

\subsection*{Blocks and Statements}

\begin{lstlisting}[language=ebnf]
block          = "{" { local_decl } { stmt } "}" ;

local_decl     = var_decl ;

stmt           = block
               | if_stmt
               | while_stmt
               | return_stmt
               | expr_stmt ;

if_stmt        = "if" "(" expr ")" stmt [ "else" stmt ] ;

while_stmt     = "while" "(" expr ")" stmt ;

return_stmt    = "return" [ expr ] ";" ;

expr_stmt      = expr ";" ;
\end{lstlisting}

\subsection*{Expressions}

\begin{lstlisting}[language=ebnf]
expr           = assignment ;

assignment     = logical_or [ "=" assignment ] ;

logical_or     = logical_and { "||" logical_and } ;

logical_and    = equality { "&&" equality } ;

equality       = relational { ("==" | "!=") relational } ;

relational     = additive { ("<" | "<=" | ">" | ">=") additive } ;

additive       = multiplicative { ("+" | "-") multiplicative } ;

multiplicative = unary { ("*" | "/" | "%") unary } ;

unary          = ("-" | "!") unary | postfix ;

postfix        = primary
                 { "(" [ arguments ] ")"
                 | "[" expr "]" } ;

arguments      = expr { "," expr } ;

primary        = IDENTIFIER
               | INT_LITERAL
               | FLOAT_LITERAL
               | CHAR_LITERAL
               | STRING_LITERAL
               | "(" expr ")" ;
\end{lstlisting}

This grammar is enough to express the required subset while remaining simple
enough for a handwritten parser.

\section{Parser Design and AST}

The parser is implemented as a recursive-descent parser. This decision was made
for three reasons:

\begin{enumerate}[leftmargin=2em]
    \item the language subset is moderate and LL-friendly after factoring
    expression precedence into separate functions,
    \item parser logic remains easy to debug and explain during presentation,
    \item explicit AST construction is useful because the project must expose
    each stage separately.
\end{enumerate}

Top-level declarations are parsed by first reading a type and identifier, then
deciding whether the declaration is a function definition or a global variable.
Blocks enforce the simplified ``declarations first, then statements'' rule.

The AST stores:

\begin{itemize}[leftmargin=2em]
    \item top-level declaration list
    \item separate lists of global variables and functions
    \item statements such as blocks, declarations, conditionals, loops, returns,
    and expression statements
    \item expressions such as literals, identifiers, calls, array indexing,
    unary operators, binary operators, and assignments
\end{itemize}

Example AST excerpt for a global array program:

\begin{lstlisting}
Program
    GlobalVarDecl name=nums type=int array=3
    FuncDecl name=main return=int
        Params <empty>
        Body
            BlockStmt
                ExprStmt
                    AssignExpr
                        Target
                            IndexExpr name=nums
                                IntLiteral value=0
\end{lstlisting}

\section{Semantic Analysis}

Semantic analysis is the stage that turns a syntactically valid program into a
semantically meaningful one. The project checks declarations, visibility,
assignment compatibility, function-call correctness, return types, and array
usage.

\subsection*{Symbol Tables and Scope}

The analyzer uses:

\begin{itemize}[leftmargin=2em]
    \item a global symbol table for global variables,
    \item a function table for user-defined functions and builtins,
    \item a stack of local scopes for nested blocks.
\end{itemize}

Lookup order is:

\begin{enumerate}[leftmargin=2em]
    \item innermost local scope,
    \item outer local scopes,
    \item global scope.
\end{enumerate}

This correctly supports shadowing. For example, a local variable named
\texttt{g} hides a global variable \texttt{g} inside the block where the local
variable is declared.

\subsection*{Checks Implemented}

The analyzer reports errors for:

\begin{itemize}[leftmargin=2em]
    \item duplicate global declarations
    \item duplicate local declarations in the same scope
    \item duplicate function declarations
    \item global/function name collisions
    \item undeclared identifiers
    \item using a non-array as an array
    \item wrong function argument count
    \item wrong argument types
    \item invalid assignment targets
    \item illegal narrowing such as \texttt{float -> int}
    \item invalid global or local initializers
    \item mismatched return values
\end{itemize}

\subsection*{Type Rules}

Safe widening conversions are allowed:

\begin{itemize}[leftmargin=2em]
    \item \texttt{char -> int}
    \item \texttt{char -> float}
    \item \texttt{int -> float}
\end{itemize}

Unsafe narrowing conversions are rejected:

\begin{itemize}[leftmargin=2em]
    \item \texttt{float -> int}
    \item \texttt{int -> char}
    \item \texttt{float -> char}
\end{itemize}

Comparison and logical expressions produce an \texttt{int} result with value
\texttt{0} or \texttt{1}. Conditions in \texttt{if} and \texttt{while} must be
numeric scalar expressions.

\subsection*{Builtins}

The following builtin signatures are pre-registered before user-defined
functions are analyzed:

\begin{lstlisting}[language=C]
print_int(int)
print_float(float)
print_char(char)
print_str(char[])
print_newline()
\end{lstlisting}

\section{Intermediate Representation}

After semantic analysis, the compiler lowers the AST to a custom intermediate
representation. This IR is already close to the VM instruction set, but still
expressed in compiler-owned data structures. The IR contains:

\begin{itemize}[leftmargin=2em]
    \item explicit global declarations
    \item explicit global initialization instructions
    \item function-level instructions
    \item jumps and labels for control flow
    \item function-call instructions
    \item local/global load and store instructions
    \item local/global array indexing instructions
\end{itemize}

For example, storing into a global array lowers into explicit indexed store
instructions rather than remaining as a high-level AST node.

\section{Object Code Generation}

The compiler executable \texttt{cforge} lowers IR into a human-readable object
format with extension \texttt{.vmo}. This object code is intentionally textual
instead of binary because it is easier to debug, diff, inspect, and present.

The object file format has three sections:

\begin{enumerate}[leftmargin=2em]
    \item \texttt{globals ... endglobals}
    \item \texttt{init ... endinit}
    \item one or more \texttt{func ... end} blocks
\end{enumerate}

Example excerpt:

\begin{lstlisting}
globals
  msg:char[6]
endglobals
init
  INIT_STRING msg "hello"
endinit
func main return=int
  params <empty>
  PUSH_GLOBAL_REF msg
  CALL_BUILTIN print_str 1
  CALL_BUILTIN print_newline 0
  PUSH_INT 0
  RET
end
\end{lstlisting}

This design gives a clear boundary between compiler and runtime, which also
helps the final demo because the generated object code can be shown directly.

\section{VM Runtime Design}

The runtime executable \texttt{crun} parses the \texttt{.vmo} object file and
executes it on a custom virtual machine.

\subsection*{Runtime Memory Model}

The VM uses:

\begin{itemize}[leftmargin=2em]
    \item a global storage map for global variables and arrays,
    \item a frame-local storage map for each function activation,
    \item a typed value stack for expression evaluation.
\end{itemize}

Each value carries a kind such as \texttt{int}, \texttt{float}, \texttt{char},
string literal, or array reference. Arrays are stored as fixed-size homogeneous
cells, and \texttt{print\_str} can accept either a string literal or a reference
to a \texttt{char[]} object.

\subsection*{Execution Flow}

The runtime performs the following steps:

\begin{enumerate}[leftmargin=2em]
    \item allocate global storage from the \texttt{globals} section,
    \item execute the \texttt{init} section once,
    \item find \texttt{main},
    \item execute \texttt{main} using local frames and the evaluation stack.
\end{enumerate}

\subsection*{Builtins in the Runtime}

Builtins are implemented inside the VM rather than being linked externally. This
keeps the system self-contained and avoids dependence on platform-specific
calling conventions.

\section{Build and Execution}

The project can be built from the repository root with:

\begin{lstlisting}
make build
\end{lstlisting}

The automated tests can be run with:

\begin{lstlisting}
make test
\end{lstlisting}

Compilation to object code:

\begin{lstlisting}
./bin/cforge samples/slice9_run.c -o /tmp/opencode/program.vmo
\end{lstlisting}

Runtime execution:

\begin{lstlisting}
./bin/crun /tmp/opencode/program.vmo
\end{lstlisting}

Phase-wise output can be obtained with:

\begin{lstlisting}
./bin/cforge --tokens  samples/slice8_run.c
./bin/cforge --ast     samples/slice8_run.c
./bin/cforge --sema    samples/slice8_run.c
./bin/cforge --ir      samples/slice8_run.c
./bin/cforge --codegen samples/slice8_run.c
\end{lstlisting}

The Makefile also provides convenience targets such as:

\begin{lstlisting}
make run FILE=samples/slice9_run.c OBJ=/tmp/opencode/program.vmo
make phases FILE=samples/slice8_run.c OBJ=/tmp/opencode/program.vmo
\end{lstlisting}

\section{Sample Programs and Outputs}

To keep the report compact while still representative, four valid programs and
one invalid program are highlighted below. The repository contains a larger pass
and fail matrix used during testing.

\subsection*{Program 1: Globals, Shadowing, and Function Use}

File: \texttt{samples/pass/global\_scalar.c}

\begin{lstlisting}[language=C]
int g = 10;

int add(int x) {
    return x + g;
}

int main() {
    print_int(add(5));
    print_newline();
    return 0;
}
\end{lstlisting}

This demonstrates:

\begin{itemize}[leftmargin=2em]
    \item global variable declaration
    \item user-defined function call
    \item global lookup inside a function body
\end{itemize}

Observed output:

\begin{lstlisting}
15
\end{lstlisting}

\subsection*{Program 2: Arrays and Strings}

File: \texttt{samples/slice8\_run.c}

\begin{lstlisting}[language=C]
int main() {
    char s[6] = "hello";
    int a[3];
    a[0] = 4;
    a[1] = 5;
    a[2] = a[0] + a[1];
    print_str(s);
    print_newline();
    print_int(a[2]);
    print_newline();
    return a[2];
}
\end{lstlisting}

This demonstrates:

\begin{itemize}[leftmargin=2em]
    \item fixed-size arrays
    \item \texttt{char[]} string initialization
    \item array indexing and assignment
    \item builtin string and integer printing
\end{itemize}

Observed output:

\begin{lstlisting}
hello
9
\end{lstlisting}

\subsection*{Program 3: Float Arithmetic and Function Return}

File: \texttt{samples/slice9\_run.c}

\begin{lstlisting}[language=C]
float addf(float a, float b) {
    return a + b;
}

int main() {
    float x = 2.5;
    float y = 4.0;
    float z = addf(x, y);
    print_float(z);
    print_newline();
    return 0;
}
\end{lstlisting}

This demonstrates:

\begin{itemize}[leftmargin=2em]
    \item float declarations and literals
    \item float-returning function
    \item builtin float printing
\end{itemize}

Observed output:

\begin{lstlisting}
6.5
\end{lstlisting}

\subsection*{Program 4: Control Flow}

File: \texttt{samples/slice5\_run.c}

\begin{lstlisting}[language=C]
int main() {
    int i = 0;
    int sum = 0;

    while (i < 5) {
        if (i == 3) {
            print_int(99);
            print_newline();
        }

        sum = sum + i;
        i = i + 1;
    }

    print_int(sum);
    print_newline();
    return sum;
}
\end{lstlisting}

This demonstrates:

\begin{itemize}[leftmargin=2em]
    \item \texttt{while} loop
    \item nested \texttt{if}
    \item jump-based lowering to the VM
\end{itemize}

Observed output:

\begin{lstlisting}
99
10
\end{lstlisting}

\subsection*{Program 5: Invalid Narrowing Conversion}

File: \texttt{samples/fail/float\_to\_int.c}

\begin{lstlisting}[language=C]
int main() {
    int x = 2.5;
    return x;
}
\end{lstlisting}

Observed compiler diagnostic:

\begin{lstlisting}
cforge: semantic: initializer type mismatch for x at 2:5: want int
\end{lstlisting}

This example demonstrates that the compiler rejects implicit narrowing from
\texttt{float} to \texttt{int}.

\section{Testing Strategy}

Testing was performed at three levels.

\subsection*{Unit Tests}

Go unit tests exist for:

\begin{itemize}[leftmargin=2em]
    \item tokenization
    \item parsing of declarations, functions, arrays, and expressions
    \item semantic analysis of declarations, calls, scopes, and type rules
    \item VM execution, jumps, function calls, object-format round trips, and
    global initialization
\end{itemize}

\subsection*{Sample Matrix}

The project contains separate \texttt{samples/pass} and \texttt{samples/fail}
directories. Valid samples confirm successful execution, while invalid samples
confirm that diagnostics are reported. The error handling is currently
\emph{fail-fast}: one run reports the first encountered semantic or syntax error.
For this reason, the negative test suite uses multiple distinct invalid sample
programs rather than a single file containing many unrelated problems.

\subsection*{Phase Verification}

Important samples were checked through:

\begin{enumerate}[leftmargin=2em]
    \item token dump
    \item AST dump
    \item semantic pass
    \item IR dump
    \item object-code dump
    \item full runtime execution
\end{enumerate}

This phase-by-phase validation helped ensure that bugs were isolated at the
correct stage rather than being observed only at the final execution step.

\section{Innovation Features}

The project satisfies the core assignment requirements and also includes
innovation-oriented work beyond the minimum requirements.

\subsection*{Core Requirements Completed}

\begin{itemize}[leftmargin=2em]
    \item lexical analysis
    \item syntax analysis
    \item semantic analysis
    \item global and local scope management
    \item intermediate representation
    \item code generation and executable runtime path
\end{itemize}

\subsection*{Innovation Features Implemented}

\begin{enumerate}[leftmargin=2em]
    \item \textbf{Implementation in Go instead of C/C++}. This directly matches
    the optional innovation category of implementing the project in another
    language.
    \item \textbf{Array support}. The compiler supports 1D fixed-size arrays of
    \texttt{int}, \texttt{char}, and \texttt{float}, including indexing and
    string initialization for \texttt{char[]}.
    \item \textbf{Custom target-code path}. Instead of stopping at IR, the
    compiler emits a standalone object-code format executed by a separate custom
    runtime. This demonstrates a complete end-to-end design beyond the minimum
    frontend-only approach.
\end{enumerate}

These additions strengthen the project in terms of completeness, modularity, and
demonstration value.

\section{Limitations and Future Work}

The current implementation is complete for the chosen subset, but several
extensions remain possible.

\begin{itemize}[leftmargin=2em]
    \item add \texttt{for} loops
    \item add 2D arrays
    \item add array parameters
    \item add an optimizer pass such as constant folding or dead code removal
    \item add escape-sequence support in char and string literals
    \item implement true short-circuit semantics for \texttt{\&\&} and \texttt{||}
    \item add global constant propagation or other global optimizations
\end{itemize}

These are natural next steps, but were intentionally deferred to keep the core
compiler stable and demonstrable by the deadline.

\section{Screenshot Placeholders}

The following placeholders can be replaced later with actual terminal
screenshots from the project directory.

\screenshotplaceholder{Compiler phase outputs for a sample program (tokens, AST, semantic analysis, IR, and object code).}{Run \texttt{make phases FILE=samples/slice8\_run.c OBJ=/tmp/opencode/program.vmo}}

\screenshotplaceholder{Compilation of a source file into a \texttt{.vmo} object file using \texttt{cforge}.}{Run \texttt{./bin/cforge samples/slice9\_run.c -o /tmp/opencode/program.vmo}}

\screenshotplaceholder{Execution of generated object code using \texttt{crun}.}{Run \texttt{./bin/crun /tmp/opencode/program.vmo}}

\screenshotplaceholder{An invalid sample and the resulting compiler diagnostic.}{Run \texttt{./bin/cforge samples/fail/float\_to\_int.c}}

\section{Conclusion}

This project delivers a full compiler pipeline for a carefully chosen C-like
subset. The final system includes lexical analysis, recursive-descent parsing,
explicit AST construction, semantic checking with both global and local scope
management, IR lowering, object-code generation, and runtime execution on a
custom VM. The use of separate compiler and runtime tools, together with array
support and a non-C/C++ implementation language, also provides meaningful
innovation beyond the core requirements. Overall, the implementation satisfies
the main project goals while remaining modular, demonstrable, and extensible.

\end{document}
